\section{Refinement workflow}
\label{sec:refinement_workflow}

Several parameter groups affect the same detector features. Detector distance changes their positions, sample alignment changes the incident condition, and mosaicity and source spread can both change widths. Refinement is staged to reduce compensation between parameter groups. Beam and detector quantities are established first, followed by sample alignment and orientation. Intensity differences are interpreted structurally after these geometric and angular constraints are established.

Table~\ref{tab:refinement_workflow} identifies the data and observables used at each stage. Instrumental parameters can be shared by measurements made in the same configuration. Sample-dependent parameters describe the film orientation, structural intensity, and, for PbI$_2$, stacking correlations.

\begin{table}[!ht]
  \centering
  \caption{Staged refinement workflow. Steps 1--5 describe the ordered Bi$_2$Se$_3$ and Bi$_2$Te$_3$ test cases. Step 6 is the added PbI$_2$ extension for stacking-fault diffuse scattering. All sample images are dark-subtracted before the listed observables are formed. A quantitative comparison requires measured and directly calculated profiles to use the same calibrated coordinate map, finite integration limits, and normalization.}
  \label{tab:refinement_workflow}
  \scriptsize
  \setlength{\tabcolsep}{3pt}
  \renewcommand{\arraystretch}{1.08}
  \begin{tabularx}{\linewidth}{@{}c >{\raggedright\arraybackslash}p{2.45cm} >{\raggedright\arraybackslash}X >{\raggedright\arraybackslash}X >{\raggedright\arraybackslash}p{2.75cm}@{}}
    \toprule
    Step & Stage & Data & Parameters refined & Primary information \\
    \midrule
    1 &
    Beam characterization &
    Direct-beam images at several detector distances &
    Beam widths $w_x$, $w_z$ and divergences $\Delta\theta_x$, $\Delta\theta_z$ &
    Incident phase space \\
    2 &
    Detector geometry calibration &
    Powder calibrant image &
    Detector distance $D$, detector tilts $(\beta_D,\gamma_D)$, and beam center $(x_0,y_0)$ &
    Detector mapping \\
    3 &
    Sample alignment and goniometer misalignment &
    Indexed reflection centers from sample images &
    Sample geometry $(\theta_i,z_B,z_S)$, in-plane rotation $\chi$, goniometer misalignment $(\alpha,\psi)$ &
    Detector-space peak positions \\
    4 &
    Mosaicity/profile refinement &
    Selected off-specular ROIs, transverse $\phi_f$ profiles, and off-condition $r=0$ regions &
    Narrow-core plus broad-component mosaic parameters &
    Off-specular arc widths, transverse $\phi_f$ widths, and off-condition axial intensity \\
    5 &
    Ordered-film comparison &
    Selected detector-space Bragg regions and $Q_z$ projections from fixed-$r$ trajectories &
    Image-wise nuisance scales and constrained ordered-intensity parameters &
    Selected Bragg-feature locations and apparent line shapes \\
    6 &
    PbI$_2$ diffuse extension &
    Fixed-$Q_R$ diffuse regions, ordered-baseline residuals, and projected $Q_z$ profiles &
    Stacking-fault probabilities or layer-pair correlation parameters &
    Diffuse intensity along $Q_z$ and separation of ordered and disorder-derived scattering \\
    \bottomrule
  \end{tabularx}
\end{table}

Direct-beam measurements constrain the beam distributions, and a powder calibrant establishes the detector mapping. Indexed sample reflections then constrain the sample pose. This geometric stage is related to detector-space indexing: it establishes where the scattering should appear before profile widths are interpreted.\cite{SmilgiesLi2026IndexGIXS}

With geometry established, off-specular and transverse axial profiles constrain the orientation distribution. Radial axial intensity depends on finite-stack scattering, orientation selection, the source distribution, and the optical response. In the model controls described in Supporting Information Sec.~S3, the radial extension persists when the source is reduced to a single monochromatic incident ray. This result does not support assigning that extension to wavelength dispersion alone.

A quantitative test of the ordered-film model requires directly calculated profiles integrated over the same finite regions as the measurements (Supporting Information Secs.~S4 and S6). For PbI$_2$, stacking variables are introduced after establishing each specimen's geometry and orientation. The added variables describe layer correlations and their effect on intensity along the rods. Source sampling is defined in Sec.~S1, finite integration in Sec.~S4, and refinement objectives, numerical convergence, and uncertainty propagation in Sec.~S8 of the Supporting Information.
